\documentclass[10pt,a4paper]{amsart}
\usepackage[margin=19mm]{geometry}
\usepackage{amsmath,amssymb,mathtools}
\usepackage[scr=rsfs,cal=euler]{mathalfa}
\usepackage{xfrac}
\usepackage[unicode,hidelinks,bookmarks=false]{hyperref}
\newcommand{\ie}{i.e.\ }
\newcommand{\cf}{cf.\ }
\newcommand{\resp}{respectively\ }
\newcommand{\Subsection}{\subsection}
\providecommand{\nobreakdash}{\nobreak\hbox{-}}
\setlength{\emergencystretch}{2em}
\multlinegap0pt
\usepackage{mathtools}
\usepackage{csquotes}
\newtheorem{theorem}{Theorem}
\usepackage{cleveref}
\crefname{theorem}{Theorem}{Theorems}
\newcommand*{\D}{\mathrm{d}} % Differential d
\newcommand*{\R}{\mathbb R}
\title{\texorpdfstring{\vspace*{-1.5\baselineskip}}{} The trace-free Einstein tensor is not variational for the metric as field variable}
\author{Arian L. von Blanckenburg, Domenico Giulini, Philip K. Schwartz}
\date{Source: arXiv:2509.02490v2 (CC BY 4.0)}
\begin{document}
\maketitle

\noindent\emph{Source and licence.} \texorpdfstring{\vspace*{-1.5\baselineskip}}{} The trace-free Einstein tensor is not variational for the metric as field variable. Version arXiv:2509.02490v2 (CC BY 4.0), distributed by arXiv under the Creative Commons Attribution 4.0 International licence, http://creativecommons.org/licenses/by/4.0/, as recorded in the arXiv licence metadata for this exact version and shown on the abstract page https://arxiv.org/abs/2509.02490v2. Full licence notice: \texttt{licenses/arxiv-2509.02490-CC-BY-4.0.txt}.\par\smallskip

\noindent\textit{Editorial scope.} Counted unit: the article's theorem thm:non-var (source0.tex lines 195--207). The article's complete original proof of it is delivered with it (source0.tex lines 325--377, one \texttt{proof} environment of 53 lines, which the article places away from the statement). No mathematics is written, repaired or re-derived: the statement and the proof are the article's own lines, in the article's own notation, and no retained line is substituted or re-flowed.  Retained uncounted as the source-local context the proof needs: lines 313--320.  Nothing was omitted from any retained span, so the package carries no omission record.  No retained line contains a citation, so the wrapper carries no bibliography and no bibliography entry.  No retained line contains a figure environment, an image-inclusion command or drawing code, so no image is delivered and the package declares no asset dependency.  Editorial formatting only, in addition: an AMS \texttt{amsart} preamble that restates the article's own declarations and macros used by the retained lines, namely the environments \texttt{theorem} on separate counters, together with the standard packages those lines need.  The retained environments print in this document as Theorem 1 in order of appearance, whereas the article prints the same items with the numbering of its own full document.  The complete native source of the article as distributed in the arXiv e-print for this exact version is bundled unchanged as \texttt{sources/184330/source0.tex}.
\begin{theorem}[informal statement]
  \label{thm:var_compl_contrap}
  If the Vainberg--Tonti Lagrangian $\mathcal L[y^A]$ associated to
  $E_A[y^B]$ exists and variation of the action $\int \mathcal L[y^A]
  \D^n x$ with respect to $y^A$ yields an expression \emph{different}
  from $E_A[y^B]$, then $E_A[y^B]$ cannot be locally variational.
  \qed
\end{theorem}

\begin{theorem}
  \label{thm:non-var}
  For $n \ge 3$, the densitised trace-free Einstein tensor
  \begin{equation}
    \label{eq:tf_Einstein_dens}
    \tilde{G}^\mathrm{tf}_{\mu\nu}
    := \sqrt{|\det(g_{\kappa\lambda})|} \, G^\mathrm{tf}_{\mu\nu}
    = \sqrt{|\det(g_{\kappa\lambda})|}
      \left(R_{\mu\nu} - \frac{1}{n} R g_{\mu\nu}\right)
  \end{equation}
  cannot arise by variation of \emph{any} local action functional with
  respect to the inverse metric.
\end{theorem}

\begin{proof}[Proof of \cref{thm:non-var}]
  We consider the densitised trace-free Einstein tensor
  \eqref{eq:tf_Einstein_dens} as an expression depending on the
  inverse metric and its derivatives, i.e.\
  $\tilde{G}^\mathrm{tf}_{\mu\nu}[g^{\rho\sigma}]$.  In order to
  compute the Vainberg--Tonti Lagrangian associated to
  $\tilde{G}^\mathrm{tf}_{\mu\nu}[g^{\rho\sigma}]$ , we have to
  determine how it behaves when the inverse metric is scaled with $u
  \in \R \setminus \{0\}$.  (Here we have to exclude the value $u = 0$
  in order to stay in the set of pseudo-Riemannian metrics.)

  By definition, when the inverse metric is scaled with $u$, the
  metric is scaled with $u^{-1}$.  Thus the Christoffel symbols are
  invariant under this scaling, likewise the Ricci tensor.  Hence the
  Ricci scalar $R = g^{\mu\nu} R_{\mu\nu}$ is scaled with $u$.
  Finally, $\sqrt{|\det(g_{\kappa\lambda})|}$ is scaled with
  $|u|^{-n/2}$.  Combined, we obtain
  \begin{equation}
    \tilde{G}^\mathrm{tf}_{\mu\nu}[u g^{\rho\sigma}]
    = |u|^{-n/2} \sqrt{|\det(g_{\kappa\lambda})|}
      \left(R_{\mu\nu} - \frac{1}{n} u R u^{-1} g_{\mu\nu}\right)
    = |u|^{-n/2} \tilde{G}^\mathrm{tf}_{\mu\nu}[g^{\rho\sigma}] \; .
  \end{equation}
  Thus, there is a unique $a \in \R \cup \{\infty\}$ satisfying
  $\lim_{u \to a} u \tilde{G}^\mathrm{tf}_{\mu\nu}[u g^{\rho\sigma}] =
  0$, namely $a = \infty$: for $n \ge 3$, the overall scaling factor
  $u |u|^{-n/2} = \mathrm{sgn}(u) |u|^{-(n-2)/2}$ vanishes only as $u
  \to \infty$.

  Hence, the Vainberg--Tonti Lagrangian associated to the densitised
  trace-free Einstein tensor as a functional of the inverse metric is
  \begin{align}
    \mathcal L[g^{\mu\nu}]
    &= g^{\rho\sigma} \int_\infty^1
      \tilde{G}^\mathrm{tf}_{\rho\sigma}[u g^{\mu\nu}] \D u \nonumber\\
    &= g^{\rho\sigma} \tilde{G}^\mathrm{tf}_{\rho\sigma}[g^{\mu\nu}]
      \int_\infty^1 u^{-n/2} \D u \nonumber\\
    &= -\frac{2}{n-2} \,
      g^{\rho\sigma} \tilde{G}^\mathrm{tf}_{\rho\sigma}[g^{\mu\nu}]
      \nonumber\\
    &= 0 \; .
  \end{align}
  In the third step, we have used that $n \ge 3$ in the computation of
  the integral, and in the last step that
  $\tilde{G}^\mathrm{tf}_{\mu\nu}$ is trace-free, i.e.\ $g^{\mu\nu}
  \tilde{G}^\mathrm{tf}_{\mu\nu} = 0$.

  Since the Vainberg--Tonti Lagrangian vanishes, variation of the
  corresponding action with respect to $g^{\mu\nu}$ yields $0$.
  Hence, according to \cref{thm:var_compl_contrap}, the densitised
  trace-free Einstein tensor is not variational for the inverse metric
  as field variable.
\end{proof}

\end{document}
